% Section 4.2, from lectures/L04/appendix/b_integrated_circuits.md.
%
% B.3 once called the noise margin "brusimmuniteten från L02"; the noise tolerance is L01 A.1, and
% the source now says so too. B.8's "nästa pass, L05" is chapter 5, where Labb 2 belongs.

\section{IC-kretsar}
\label{c4:app:b}

\subsection{Från grind till krets}
\label{c4:b1}
En logisk grind är ett begrepp. För att kunna koppla upp den behöver du en \textbf{IC-krets}
(\emph{integrated circuit}): en liten kiselbit med transistorer, ingjuten i en plastkapsel med
metallben. Den mest spridda familjen av enkla logikkretsar är \textbf{74-serien}, som funnits sedan
1960-talet och fortfarande tillverkas. Varje krets i serien har ett nummer som säger vad den
innehåller:

\begin{booktable}{@{}p{5.2em}Lp{6em}@{}}
  \thead{Krets} & \thead{Innehåll} & \thead{Används i} \\
  \midrule
  74HC00 & fyra NAND-grindar med två ingångar & Labb~2-2 \\
  74HC02 & fyra NOR-grindar med två ingångar & Labb~2-2 \\
  74HC04 & sex inverterare (NOT) & Labb~2-1 \\
  74HC08 & fyra AND-grindar med två ingångar & Labb~2-1 \\
  74HC32 & fyra OR-grindar med två ingångar & Labb~2-1 \\
  74HC86 & fyra XOR-grindar med två ingångar & -- \\
\end{booktable}

Namnet har tre delar. \textbf{74} betyder standardserien, för normala temperaturer. \textbf{HC} är
kretsfamiljen, alltså hur transistorerna inuti är byggda: \emph{High-speed CMOS}. \textbf{00} är
funktionen. En 74HC00 och en 74LS00 innehåller samma fyra NAND-grindar, men i olika teknik.

Tillverkarna lägger ofta till egna bokstäver före och efter. \code{SN74HC00N} från Texas
Instruments och \code{74HC00N} från Nexperia är samma krets; \code{N} på slutet betyder här kapseln
DIP, den som passar i ett kopplingsdäck. Leta efter \code{74}, familjen och numret, och bry dig inte
om resten.

\textbf{Två familjer du kommer att möta:}
\begin{itemize}
  \item \textbf{HC} (CMOS) arbetar på 2--6~V, drar nästan ingen ström i vila, och är den boken
    använder.
  \item \textbf{LS} (TTL) är äldre, kräver 5~V, och drar mer ström. En LS-ingång som lämnas
    okopplad läses som 1, vilket har lärt generationer av elektroniker den dåliga vanan att lämna
    ingångar okopplade. På en HC-krets är det ett fel (\secref{c4:b4}).
\end{itemize}

\subsection{Kapsel och benplacering}
\label{c4:b2}
Kretsarna i boken sitter i en \textbf{DIP-14}-kapsel: 14 ben i två rader, 2,54~mm mellan benen, så
att de passar i kopplingsdäckets hål.

\textbf{Ben 1 hittar du med hjälp av markeringen.} I ena kortsidan finns ett halvrunt
\textbf{hack}, och ofta också en liten prick vid ben 1. Håll kretsen med hacket åt vänster och
texten rättvänd: då sitter ben 1 längst ned till vänster. Benen räknas sedan \textbf{moturs, sett
ovanifrån}: 1 till 7 längs nederkanten från vänster till höger, och 8 till 14 längs överkanten från
höger till vänster.

Alla kretsar i tabellen ovan har matningen på samma ben:
\begin{itemize}
  \item \textbf{ben 14: VCC}, plus 5~V;
  \item \textbf{ben 7: GND}, jord (0~V).
\end{itemize}

Utan matning fungerar ingenting. Och en krets som sitter felvänd får plus på ben 7 och jord på ben
14: den blir varm, fungerar inte, och kan förstöras. Känn efter med ett finger om en krets blir
varm, och bryt matningen direkt om den blir det.

Benplaceringen finns i databladet. Här är de kretsar boken använder, ritade ovanifrån med grindarna
inuti:

\bookfigure{l04_dip_74hc00}{Benplacering för 74HC00, fyra NAND-grindar.}{c4:fig:dip_74hc00}

\bookfigure{l04_dip_74hc08}{Benplacering för 74HC08, fyra AND-grindar: samma benplacering som
74HC00.}{c4:fig:dip_74hc08}

\bookfigure{l04_dip_74hc32}{Benplacering för 74HC32, fyra OR-grindar: samma benplacering som
74HC00.}{c4:fig:dip_74hc32}

\bookfigure{l04_dip_74hc86}{Benplacering för 74HC86, fyra XOR-grindar: samma benplacering som
74HC00.}{c4:fig:dip_74hc86}

\bookfigure{l04_dip_74hc04}{Benplacering för 74HC04, sex inverterare.}{c4:fig:dip_74hc04}

\bookfigure{l04_dip_74hc02}{Benplacering för 74HC02, fyra NOR-grindar: utgången kommer först i
varje grupp.}{c4:fig:dip_74hc02}

\textbf{74HC00, 08, 32 och 86 har samma benplacering.} Grind 1 har ingångarna på ben 1 och 2 och
utgången på ben 3, grind 2 på ben 4, 5 och 6, och så vidare. Byt en 74HC08 mot en 74HC32 så byts
alla fyra AND-grindarna mot OR-grindar, utan att en enda sladd flyttas.

\textbf{74HC02 är undantaget}, och det är den krets flest kopplar fel. I varje grupp om tre ben
kommer \textbf{utgången först}: ben 1 är utgång och ben 2 och 3 ingångar. Den som kopplar en 74HC02
ur minnet av en 74HC00 kopplar sina insignaler till en utgång. Titta på benplaceringen varje gång.

\subsection{Databladet}
\label{c4:b3}
Varje krets har ett \textbf{datablad}, ett dokument från tillverkaren som beskriver den i detalj.
Sök på kretsens namn och \enquote{datasheet}, och välj en stor tillverkare som Texas Instruments,
Nexperia eller onsemi. Ett datablad är ofta 15--20 sidor, men för att koppla upp en krets behöver
du bara tre saker ur det:
\begin{enumerate}
  \item \textbf{Benplaceringen} (\emph{pin configuration} eller \emph{pinout}): vilket ben som är
    vad, som i figurerna ovan.
  \item \textbf{Funktionstabellen} (\emph{function table}): kretsens sanningstabell, skriven med
    \code{H} för hög och \code{L} för låg, och ibland \code{X} för \enquote{spelar ingen roll}.
  \item \textbf{De elektriska nivåerna}: vilka spänningar som räknas som 0 och 1.
\end{enumerate}

\subsubsection{Logiknivåer}
En digital krets delar in spänningen i tre områden: låg, hög, och ett förbjudet område däremellan.
För en 74HC-krets matad med 4,5~V anger databladet:

\begin{booktable}{@{}p{4.2em}Lp{10.2em}@{}}
  \thead{Storhet} & \thead{Betydelse} & \thead{Värde vid VCC = 4,5~V} \\
  \midrule
  $V_{\text{IH}}$ & lägsta spänning på en ingång som garanterat läses som 1 & 3,15~V \\
  $V_{\text{IL}}$ & högsta spänning på en ingång som garanterat läses som 0 & 1,35~V \\
  $V_{\text{OH}}$ & lägsta spänning en hög utgång ger, vid 4~mA last & 3,84~V \\
  $V_{\text{OL}}$ & högsta spänning en låg utgång ger, vid 4~mA last & 0,33~V \\
\end{booktable}

En tumregel för HC-kretsar: under \textbf{30~\%} av matningen är 0, över \textbf{70~\%} är 1. Vid
5~V betyder det under 1,5~V och över 3,5~V. En ingång mellan de två gränserna kan läsas som vad som
helst, och får kretsen att dra onödigt mycket ström.

Lägg märke till marginalen: en utgång ger minst 3,84~V när den är hög, och en ingång behöver bara
3,15~V. Skillnaden, ungefär 0,7~V, är \textbf{brusmarginalen}. Den gör att en störning på upp till
0,7~V på en ledning inte ändrar något. Det är brusimmuniteten från \secref{c1:a1}, uttryckt i volt.

\subsubsection{Gränsvärden}
Databladet har också en tabell med \textbf{absoluta maxvärden} (\emph{absolute maximum ratings}). De
är inte arbetsvärden, utan gränser som förstör kretsen om de överskrids. För 74HC:
\begin{itemize}
  \item matningen får vara högst 7~V (och kretsen är specificerad för 2--6~V);
  \item en utgång får leverera högst 25~mA;
  \item hela kretsen, alla utgångar tillsammans, högst 50~mA genom VCC eller GND.
\end{itemize}

\textbf{Koppla aldrig 24~V från Labb~1 till en 74HC-krets.} Den överlever det inte.

\subsection{Ingångar: aldrig flytande}
\label{c4:b4}
En CMOS-ingång drar nästan ingen ström, under en mikroampere. Det låter bra, men det betyder att en
ingång som inte är ansluten till något, en \textbf{flytande ingång}, inte har någon bestämd
spänning alls. Den plockar upp laddning från omgivningen och kan läsas som 0 i ena sekunden och 1 i
nästa, beroende på om din hand är nära. Kretsen verkar fungera ibland, och det är den värsta sortens
fel.

\textbf{Regeln: varje ingång ska vara ansluten till något.} Det gäller också ingångarna på grindar
du inte använder: koppla dem till GND (eller VCC).

En strömbrytare eller tryckknapp ger bara en förbindelse när den är sluten. När den är öppen är
ingången flytande, om inte ett motstånd bestämmer nivån:

\bookfigure{l04_io_circuit}{Tryckknapp med pull-down-motstånd på en grindingång, och lysdiod med
förkopplingsmotstånd på utgången.}{c4:fig:io_circuit}

\begin{itemize}
  \item Knappen sitter mellan \textbf{5~V} och ingången.
  \item Ett motstånd på \textbf{10~kΩ} sitter mellan ingången och \textbf{jord}. Det kallas ett
    \textbf{pull-down-motstånd}: det drar ned ingången till 0 när knappen är öppen.
  \item När knappen trycks ned kopplar den ingången direkt till 5~V. Motståndet ligger då mellan
    5~V och jord och leder 0,5~mA, men det spelar ingen roll; ingången ser 5~V och läser 1.
\end{itemize}

Det omvända, knappen mot jord och motståndet mot 5~V, kallas ett \textbf{pull-up-motstånd} och ger
omvänd logik: nedtryckt knapp läses som 0. Det är den kopplingen mikrokontrollern använder i
kapitel~\ref{c10:ch}, där motståndet dessutom sitter inbyggt i kretsen.

Kopplingsdäckets stationer i Labb~2 har strömbrytare som redan ger 0 eller 5~V. Då behövs inga egna
motstånd, men regeln om ingångar du inte använder gäller fortfarande.

\subsection{Utgångar: lysdiod med förkopplingsmotstånd}
\label{c4:b5}
För att se vad en utgång gör kopplar du en \textbf{lysdiod} till den. En lysdiod är en diod som
lyser när ström går igenom den i rätt riktning: från det långa benet (\textbf{anoden}) till det
korta (\textbf{katoden}, som också är märkt med en platt kant på kapseln).

En lysdiod måste alltid ha ett \textbf{förkopplingsmotstånd} i serie. Över en röd lysdiod ligger
ungefär 2~V när den lyser, nästan oberoende av strömmen, så utan motstånd skulle strömmen bara
begränsas av hur mycket utgången orkar leverera, och det är mer än lysdioden tål.

Motståndet räknas ut med Ohms lag. Resten av spänningen, det som inte ligger över lysdioden, ska
ligga över motståndet:
\[
  R = \frac{U_{\text{matning}} - U_{\text{lysdiod}}}{I}
\]
Med 5~V matning, en röd lysdiod på 2~V och en önskad ström på 10~mA:
\[
  R = \frac{5 - 2}{0,010} = 300\ \Omega
\]
Välj närmaste standardvärde uppåt, \textbf{330~Ω}, som ger 9~mA. En lysdiod lyser tydligt redan vid
några milliampere, så 330--470~Ω är ett bra val till en 74HC-utgång. Håll strömmen under 10~mA: då
ligger utgången väl inom sina gränser, och en hög utgång håller sin spänning.

Lysdioden kopplas från utgången, genom motståndet, till jord. Då lyser den när utgången är
\textbf{1}. Vänd den, eller koppla den mot 5~V i stället för mot jord, så lyser den när utgången är
0: det fungerar också, men gör varje avläsning baklänges.

\subsection{Kopplingsdäcket}
\label{c4:b6}
Ett \textbf{kopplingsdäck} (\emph{breadboard}) låter dig koppla upp en krets utan att löda. Under
hålen finns metallfjädrar som förbinder hålen i ett bestämt mönster:

\bookfigure{l04_breadboard}{Kopplingsdäck med matningsskenor, femhålskolumner och en IC-krets över
mittspåret.}{c4:fig:breadboard}

\begin{itemize}
  \item \textbf{Matningsskenorna} längs kanterna, märkta \code{+} (röd linje) och \code{-} (blå
    linje), är förbundna längs hela kanten. Anslut 5~V och jord till dem en gång, så har du
    matning överallt.
  \item \textbf{Hålen i mitten} är förbundna \textbf{fem och fem i kolumner}, a--e på ena sidan av
    mittspåret och f--j på den andra. Allt som sitter i samma kolumn på samma sida är kopplat
    ihop.
  \item \textbf{Mittspåret} skiljer de två halvorna åt. En IC-krets sätts \textbf{över} mittspåret,
    så att varje ben hamnar i en egen kolumn, och du får fyra lediga hål bredvid varje ben att
    koppla till.
\end{itemize}

Tre vanliga fel, som alla ger en krets som \enquote{nästan} fungerar:
\begin{itemize}
  \item Två ledningar i \textbf{samma kolumn} som inte skulle vara ihopkopplade. De är det nu.
  \item En krets som \textbf{inte} sitter över mittspåret: då kortsluts benen parvis.
  \item Matningsskenor som är \textbf{delade på mitten}. På en del kopplingsdäck går skenan inte
    hela vägen, vilket syns som ett avbrott i den röda eller blå linjen. Koppla i så fall ihop
    halvorna.
\end{itemize}

Använd gärna \textbf{röd tråd för 5~V och svart eller blå för jord}, och korta trådar som ligger
platt mot däcket. En uppkoppling som går att följa med blicken går också att felsöka.

Sätt en \textbf{avkopplingskondensator} på 100~nF mellan VCC och GND, så nära varje krets som
möjligt. När en grind slår om drar den en kort strömpuls, och kondensatorn levererar den lokalt, i
stället för att pulsen ska störa matningen till resten av kretsarna. I små uppkopplingar som i
Labb~2 fungerar det ofta utan, men det är god vana, och i en verklig konstruktion är den inte
valfri.

\subsection{Att planera en uppkoppling}
\label{c4:b7}
Den som börjar koppla direkt från ett grindnät tappar bort sig efter tredje sladden. Planera i
stället i tre steg, med papper och penna.

\textbf{1. Räkna grindarna, välj kretsar.} Räkna hur många grindar av varje sort nätet behöver, och
hur många kretsar det blir. Varje 74HC08, 32, 00, 02 och 86 har fyra grindar, 74HC04 har sex.

\textbf{2. Gör en grindtilldelning.} Bestäm vilken grind i vilken krets som gör vad, och skriv ut
benens nummer. Kretsarna får beteckningarna U1, U2 och så vidare. För \code{X = AB + C}:

\begin{booktable}{@{}>{\raggedright\arraybackslash}p{6.4em}p{6.6em}%
    >{\raggedright\arraybackslash}p{4.6em}L>{\raggedright\arraybackslash}p{4.4em}@{}}
  \thead{Grind i nätet} & \thead{Krets} & \thead{Grind i kretsen} & \thead{Ingångar (ben)} &
    \thead{Utgång (ben)} \\
  \midrule
  AND: \code{AB} & U1 = 74HC08 & grind 1 & A $\rightarrow$ 1, B $\rightarrow$ 2 & 3 \\
  OR: \code{AB + C} & U2 = 74HC32 & grind 1 & U1 ben 3 $\rightarrow$ 1, C $\rightarrow$ 2 &
    3 $\rightarrow$ X \\
\end{booktable}

\textbf{3. Gör en kopplingslista.} En rad per ledning, i den ordning du kopplar. Börja alltid med
matningen:

\begin{booktable}{@{}p{2em}Lp{8em}@{}}
  \thead{Nr} & \thead{Från} & \thead{Till} \\
  \midrule
  1 & U1 ben 14 & +5~V \\
  2 & U1 ben 7 & GND \\
  3 & U2 ben 14 & +5~V \\
  4 & U2 ben 7 & GND \\
  5 & strömbrytare A & U1 ben 1 \\
  6 & strömbrytare B & U1 ben 2 \\
  7 & U1 ben 3 & U2 ben 1 \\
  8 & strömbrytare C & U2 ben 2 \\
  9 & U2 ben 3 & lysdiod X \\
  10 & oanvända ingångar (U1 ben 4, 5, 9, 10, 12, 13; U2 ben 4, 5, 9, 10, 12, 13) & GND \\
\end{booktable}

Bocka av varje rad när ledningen sitter. Om något inte fungerar har du nu en lista att kontrollera
rad för rad, och det är precis vad felsökningen i \secref{c5:app:a} bygger på.

\subsection{Förbered Labb~2}
\label{c4:b8}
Labb~2 (bilaga~\ref{app:lab2}) hör till kapitel~\ref{c5:ch}, och du hinner inte klart om
förberedelserna görs vid labbplatsen. Före labben ska du ha:
\begin{itemize}
  \item läst hela labbhandledningen, och tagit fram sanningstabeller, Karnaughdiagram och
    minimerade uttryck för varje uppgift;
  \item ritat grindnätet för varje uppgift, och byggt och testat det i CircuitVerse;
  \item gjort en grindtilldelning och en kopplingslista (\secref{c4:b7}) för varje uppgift;
  \item för Labb~2-2: gjort om näten till enbart NAND (\secref{c4:a9}), och gjort
    grindtilldelningen för 74HC00.
\end{itemize}

Ta med allt på papper eller i datorn. Förberedelserna är en del av redovisningen.
